\subsection{Component-Wise Stochastic Advection Kernel}
\label{subsec:component_advection_kernel}

The key quantity in the Lagrangian IDE kernel is the spatial displacement
between a source location and a target location. Let
\[
\mathbf h_{ij}=s_i^\ast-s_j^\ast\in\mathbb R^2
\]
denote the two-dimensional displacement from source location \(s_j^\ast\) to
target location \(s_i^\ast\). Since this displacement is defined in physical
space, the advection term that shifts the kernel must also be a
two-dimensional spatial displacement. We therefore model advection at the level
of the two component fields rather than as four scalar component-pair
displacements.

Let \(u\) and \(v\) denote the zonal and meridional wind-component fields. At
time \(t\), VectorWIDE represents the stochastic advection of these two fields
by
\begin{equation}
\boldsymbol a_t
=
\left[
\left(\boldsymbol a_t^{(u)}\right)^\top,
\left(\boldsymbol a_t^{(v)}\right)^\top
\right]^\top
\in\mathbb R^4,
\qquad
\boldsymbol a_t^{(u)},\boldsymbol a_t^{(v)}\in\mathbb R^2 .
\end{equation}
The four-dimensional vector is thus the stacking of two two-dimensional
spatial advections:
\[
\boldsymbol a_t =
(a_{u,x,t},a_{u,y,t},a_{v,x,t},a_{v,y,t})^\top .
\]
We assume
\begin{equation}
\boldsymbol a_t
\sim
\mathcal N(\boldsymbol\mu_t,\boldsymbol\Sigma_t),
\qquad
\boldsymbol\mu_t\in\mathbb R^4,\quad
\boldsymbol\Sigma_t\in\mathbb R^{4\times4}.
\label{eq:component_advection_distribution}
\end{equation}
The covariance matrix \(\boldsymbol\Sigma_t\) captures uncertainty within each
component-field advection and dependence between the two component-field
advections.

Define the component selectors
\[
\mathbf E_u =
\begin{bmatrix}
1 & 0 & 0 & 0\\
0 & 1 & 0 & 0
\end{bmatrix},
\qquad
\mathbf E_v =
\begin{bmatrix}
0 & 0 & 1 & 0\\
0 & 0 & 0 & 1
\end{bmatrix}.
\]
For a transition from source component \(q\) to target component \(p\), the
stacked advection vector is projected to physical two-dimensional space by a
matrix \(\mathbf P_{pq}\in\mathbb R^{2\times4}\). Following the Lagrangian
cross-covariance construction, the diagonal and off-diagonal blocks use
different time projections. For within-component propagation,
\begin{equation}
\mathbf P_{pp}(u)=u\mathbf E_p,
\end{equation}
where \(u\) is the time lag. For cross-component propagation,
\begin{equation}
\mathbf P_{pq}(t_1,t_2)=t_1\mathbf E_p-t_2\mathbf E_q,
\qquad p\neq q .
\label{eq:component_cross_projection}
\end{equation}
This is the analogue of the \(T(t_1,t_2)\widetilde{\mathbf d}_{ij}\)
projection in the multivariate Lagrangian theorem: the target component
advection and source component advection both contribute to the cross block.
The same projection must be used for the mean displacement and for the
covariance-induced dispersion.

Using \(\mathbf P_{pq}\), define
\begin{equation}
\mathbf h_{ij,t}^{\prime(pq)}
=
\mathbf h_{ij}
-
\mathbf P_{pq}\boldsymbol\mu_t .
\label{eq:component_pairwise_hprime}
\end{equation}
The effective dispersion matrix is
\begin{equation}
\mathbf D_t^{(pq)}
=
\ell_{pq}^2\mathbf I_2
+
2\mathbf P_{pq}\boldsymbol\Sigma_t\mathbf P_{pq}^{\top},
\label{eq:component_effective_dispersion}
\end{equation}
where \(\ell_{pq}>0\) is a component-pair baseline spatial spread. In the
one-step implementation, \(u=\Delta t\), \(t_1=\Delta t\), and
\(t_2=0\), so the cross block is shifted by the target component's effective
advection. More generally, retaining \(t_2>0\) gives the full
cross-component projection \(t_1\mathbf E_p-t_2\mathbf E_q\). The same
projection appears in both \(\mathbf h_{ij,t}^{\prime(pq)}\) and
\(\mathbf D_t^{(pq)}\). The discretized redistribution kernel is
\begin{equation}
\widetilde{k}_t^{(pq)}(s_i^\ast,s_j^\ast)
=
\alpha_t^{(pq)}
\left|\mathbf D_t^{(pq)}\right|^{-1/2}
\exp
\left\{
-
\left(\mathbf h_{ij,t}^{\prime(pq)}\right)^\top
\left(\mathbf D_t^{(pq)}\right)^{-1}
\mathbf h_{ij,t}^{\prime(pq)}
\right\},
\label{eq:component_pairwise_kernel}
\end{equation}
where \(\alpha_t^{(pq)}\) is a learned component-mixing weight. The diagonal
blocks describe within-component propagation, while the off-diagonal blocks
describe cross-component transfer. In all cases, the geometry of the spatial
kernel is governed by a two-dimensional advection displacement.

Evaluating \eqref{eq:component_pairwise_kernel} over all source-target pairs
gives the four unnormalized blocks
\[
\widetilde{\mathbf K}_t^{(pq)}
=
\left[
\widetilde{k}_t^{(pq)}(s_i^\ast,s_j^\ast)
\right]_{i,j=1}^n,
\qquad p,q\in\{u,v\}.
\]
The full transition matrix is assembled as
\begin{equation}
\widetilde{\mathbf K}_t
=
\begin{bmatrix}
\widetilde{\mathbf K}_t^{(uu)} & \widetilde{\mathbf K}_t^{(uv)}\\
\widetilde{\mathbf K}_t^{(vu)} & \widetilde{\mathbf K}_t^{(vv)}
\end{bmatrix}.
\end{equation}
We then row-normalize the full block matrix:
\begin{equation}
\left[\mathbf K_t\right]_{ij}
=
\frac{
\left[\widetilde{\mathbf K}_t\right]_{ij}
}{
\sum_{j'=1}^{2n}
\left[\widetilde{\mathbf K}_t\right]_{ij'}
},
\qquad i,j=1,\ldots,2n.
\end{equation}
This normalization yields a stable redistribution operator while preserving the
relative spatial asymmetry and component-mixing structure induced by the
learned stochastic advection moments.

\subsection{Learning Component Advection from Exogenous Wind Fields}
\label{subsec:component_advection_learning}

Let \(\mathbf X_t\) denote the exogenous NWP wind-field input available at time
\(t\). VectorWIDE uses a neural mapping to convert recent NWP fields into the
low-dimensional stochastic advection moments that parameterize the IDE kernel:
\begin{equation}
(\boldsymbol\mu_t,\boldsymbol\Sigma_t,\mathbf A_t)
=
\psi(\mathbf X_{1:t};\boldsymbol\Phi),
\label{eq:component_deep_mapping}
\end{equation}
where \(\mathbf A_t=\{\alpha_t^{(pq)}\}_{p,q\in\{u,v\}}\) contains the
component-mixing weights. In the implementation, a CNN first encodes the
spatial structure of each NWP field and a causal transformer then models the
temporal evolution of the encoded features. The output heads produce the
component-wise mean advection \(\boldsymbol\mu_t\), a positive-definite
covariance matrix \(\boldsymbol\Sigma_t\), and the component-mixing weights
\(\mathbf A_t\).

The neural network therefore does not directly replace the state-space wind
forecast. Instead, it supplies physically interpretable, time-varying
advection moments to the probabilistic IDE transition operator. This separation
keeps the transport geometry explicit: advection shifts the spatial field in
two-dimensional physical space, while the multi-output IDE structure models
within-component propagation and cross-component exchange.
